\section{Conclusion}
% 本文针对 MXFP4 作为下一代低精度浮点格式在 LLM 激活量化中面临的核心挑战，提出了 TORQ 两级正交旋转预处理框架。通过系统化分析揭示跨块方差不均衡与块内 FP4 码字使用不均的结构性问题，我们创新性地将量化优化从局部参数调节提升至坐标系设计层面：块间旋转基于 Schur-Horn 定理与凸优化理论，实现跨块方差的全局均衡分配，消除少数高方差块对总误差的支配；块内旋转通过 log₂ 域目标引导与线性域正交变换的协同设计，使激活对数幅值分布逼近均匀，最大化 MXFP4 码字利用率。理论上，我们为两级旋转分别建立了误差下界与优化性质的严格证明；实践中，该方法无需微调、校准成本低且硬件友好，在 LLaMA、Qwen等系列模型及 WikiText、MMLU、GSM8K 等多类任务上，显著缩小了 MXFP4 量化与全精度推理的精度差距，性能超越现有 SOTA 方法 2–5 %。TORQ 不仅为 MXFP4 激活量化提供了首个系统性解决方案，更构建了 “结构感知的正交变换优化” 新范式，为低比特浮点格式在 LLM 高效部署中的应用奠定了关键基础。
This paper addresses the core challenges faced by MXFP4, as a next-generation low-precision floating-point format, in the activation quantization of Large Language Models (LLMs) and proposes the TORQ two-level orthogonal rotation preprocessing framework. Through systematic analysis, it reveals the structural issues of uneven cross-block variance and unequal usage of intra-block FP4 codewords. We innovatively elevate quantization optimization from the level of local parameter adjustment to that of coordinate system design: Inter-block rotation, based on the Schur-Horn theorem and convex optimization theory, achieves the global balanced distribution of cross-block variance, eliminating the dominance of a few high-variance blocks over the total error; Intra-block rotation, through the collaborative design of $log2$ domain target guidance and linear domain orthogonal transformation, makes the distribution of activation logarithmic amplitudes approximate uniformity, maximizing the utilization of MXFP4 codewords. Theoretically, we establish strict proofs for the error lower bounds and optimization properties of the two-level rotations respectively; In practice, this method requires no fine-tuning, has low calibration costs, and is hardware-friendly. On a series of models such as LLaMA and Qwen, as well as various tasks including WikiText, MMLU, and GSM8K, it significantly narrows the precision gap between MXFP4 quantization and full-precision inference, with performance surpassing existing SOTA methods by 2–5\%. TORQ not only provides the first systematic solution for MXFP4 activation quantization but also constructs a new paradigm of "structure-aware orthogonal transformation optimization," laying a crucial foundation for the application of low-bit floating-point formats in the efficient deployment of LLMs.
